\section*{Appendix. Quasi-geodesic semi-groups}
\refstepcounter{section}
\subsection{Goal}

If \(\Gamma\) is the free group with generators \(a,b\) then the semi-group \(S\) consisting of reduced words which begin and end with \(a\) has the following nice properties:
\begin{enumerate}
\item \(S\) generates \(\Gamma\) as a group.
\item For every \(n\) and \(\gamma_1,\ldots, \gamma_n \in S\) the geodesic joining the neutral element \(e\) to \(\gamma_1\cdots \gamma_n\) passes through all partial products \(\gamma_1\cdots \gamma_k\) for \(k = 1,\ldots, n\).
\item For every \(\gamma \in \Gamma\) there exist two letter words \(L(\gamma)\) and \(R(\gamma)\) such that \(L(\gamma)\gamma R(\gamma) \in S\).
\end{enumerate}
We may for instance set \(L(\gamma)\) to be \(a\) if the reduced form of \(\gamma\) does not begin with~\(a^{-1}\) and \(ab\) otherwise, and \(R(\gamma)\) similarly to be \(a\) if the last letter of \(\gamma\) is not \(a^{-1}\) and \(ba\) otherwise.

Our goal in this appendix is to give a similar construction for a general hyperbolic group.
For this purpose we fix a group \(\Gamma\) with a finite symmetric generating set \(\cS\). We~let \(|\gamma|\) be the word length of \(\gamma \in \Gamma\) with respect to \(\cS\) and consider the word distance \(\dist(\gamma,\eta) = |\gamma^{-1}\eta|\). We~assume that there exists \(\delta > 0\) (fixed from now on) such that \(\Gamma\) is \(\delta\)-hyperbolic with respect to \(\dist\), and non-elementary (\ie contains two independent loxodromic elements).

Recall that given positive constants \(C,D > 0\) a \((C,D)\)-quasi-geodesic is a sequence \(x: I \cap \Z \to \Gamma\) where \(I\) is an interval, satisfying
\[
C^{-1}|m-n| - D \le \dist(x_m,x_n) \le C|m-n| + D,
\]
for all \(m,n \in I \cap \Z\).

We say a semigroup \(S \subset \Gamma\) is quasi-geodesic if there exist constants \(C,D > 0\) such that for all \(n\) and \(\gamma_1,\ldots, \gamma_n \in S\) the sequence
\[
e, \gamma_1,\gamma_1\gamma_2,\ldots, \gamma_1\cdots \gamma_n,
\]
is visited in left to right order by some \((C,D)\)-quasi-geodesic.

We now state the main result of the appendix.

\begin{proposition}\label{semigrouplemma}
Let \((\Gamma,\cS)\) be a non-elementary hyperbolic group with a finite symmetric generating set \(\cS\). Then there exists a semigroup \(S \subset \Gamma\), a finite set \(F \subset \Gamma\) and mappings \(L,R:\Gamma \to F\) with the following properties:
\begin{enumerate}
\item \(S\) generates \(\Gamma\) as a group,
\item \(S\) is quasi-geodesic, and
\item \(L(\gamma)\gamma R(\gamma) \in S\) for all \(\gamma \in \Gamma\).
\end{enumerate}
\end{proposition}

The arguments in what follows use the local-to-global principle for hyperbolic groups as in \cite[\S 3]{gouezel}.

\subsection{Proof of Proposition \ref{semigrouplemma}}

Recall that the Gromov product (based at the neutral element \(e \in \Gamma\)) is defined by
\[
(\gamma,\eta) = \frac{|\gamma| + |\eta| - |\gamma^{-1}\eta|}{2}.
\]
We assume \(\delta > 0\) is such that
\[
\min\lbrace (\gamma_1,\gamma_2),(\gamma_2,\gamma_3)\rbrace \le (\gamma_1,\gamma_3) + \delta,
\]
for all \(\gamma_1,\gamma_2,\gamma_3 \in \Gamma\). We~further assume that
\[
\dist(e,\alpha) \le (p,q) + \delta,
\]
for all \(p,q \in \Gamma\) and geodesic \(\alpha\) joining them (such a choice of \(\delta\) is possible by \cite[Prop.\,1.22]{bridson-haefliger}).

\begin{lemma}
For each \(C \ge 1\) there exist \(a,b \in \Gamma\) such that setting \(A = \lbrace a,b,a^{-1},b^{-1}\rbrace\) and defining
\[
c_1 = \max_{x,y \in A:\ x \neq y}(x,y)\quand c_2 = \min_{x \in A}|x|,
\]
one has \(c_1 + 2\delta + r < \frac{1}{2}{c_2}\) for some \(r > C(c_1+2\delta+1)\).
\end{lemma}
\begin{proof}
Since \(\Gamma\) is non-elementary we may take \(x,y\) two independent loxodromic elements (\ie the set of boundary fixed points of \(x\) and \(y\) are disjoint).
Letting \(a = x^n\) and \(b = y^n\) for sufficiently large \(n\) establishes the claim since \(c_1(x^n,y^n)\) remains bounded, while \(c_2(x^n,y^n)\) goes to infinity.
\end{proof}

We fix from now on a set \(A\) as in the previous lemma with corresponding constants \(c_1,c_2,r\) satisfying
\begin{equation}\label{requation}
r > 32(c_1+2\delta+1).
\end{equation}

\begin{lemma}\label{borrowedlemma}
The sets defined for \(x \in A\) by \(V(x) = \lbrace \gamma: (\gamma,x) \ge c_1+2\delta + r\rbrace\) are pairwise disjoint.

\begin{itemize}
\item
For all \(x \in A\) and \(\gamma \notin V(x^{-1})\) one has \(x\gamma \in V(x)\).

\item
For all \(x,y \in A\) with \(x \neq y\), if~\(\gamma \in V(x)\) and \(\eta \in V(y)\) then \((\gamma,\eta) \le c_1+2\delta\).

\item
If \(\gamma \in V(x)\) for some \(x \in A\) then \(|\gamma| \ge r\).
\end{itemize}
\end{lemma}
\begin{proof}
By hypothesis \(2\delta + r \in (2\delta,c_2/2 - c_1)\). With this property, the first three statements are claims 1--3 in the proof of \cite[Lem.\,4.1]{random}.

For the last statement observe that
\[
|\gamma| = (\gamma,\gamma) \ge \min\lbrace (\gamma,x),(x,x)\rbrace - \delta \ge \min \lbrace c_1 + 2\delta + r, c_2\rbrace - \delta \ge r.\qedhere
\]
\end{proof}

\begin{lemma}\label{mlemma}
There exists \(m\) such that for all \(\gamma \!\in\! V(a)\) and all \(\eta\) with \(\dist(\gamma,\eta) \!\le\! 2c_2\) one has \(a^m\eta \in V(a)\).
\end{lemma}
\begin{proof}
We first observe that by Lemma \ref{borrowedlemma} we have \((a^{-1},a) \le c_1 + 2\delta\) and therefore a geodesic fixed under multiplication by \(a\) is at distance at most \(c_1+3\delta\) from \(a^m\) for all \(m\).
It follows, letting \(\ell = \lim_{n \to +\infty}\frac{1}{n}|a^n|\) be the translation length of \(a\), that
\[
\ell \ge |a| - 2(c_1+3\delta) \ge 2(c_1+2\delta+r) -2(c_1+3\delta) \ge 2r-2\delta.
\]
Hence we obtain for all \(m \ge 1\) that
\begin{align*}
(a,a^m) &= \frac{|a|+|a^m|-|a^{m-1}|}{2}
\\ &\ge \frac{\ell -2(c_1+3\delta) + m\ell - 2(c_1+3\delta) - (m-1)\ell -2(c_1+3\delta)}{2}
\\ &= \ell - 3(c_1+3\delta)
\ge 2r - 3c_1-11\delta
\ge c_1+3\delta + r.
\end{align*}

In what follows we will need to switch basepoints, so we recall that the Gromov product
\[
(x,y)_z = \frac{\dist(x,z)+\dist(y,z)-\dist(x,y)}{2},
\]
satisfies
\((x,y)_z + (z,y)_x = \dist(x,z)\).
We now calculate using Lemma \ref{borrowedlemma}
\begin{align*}
(a^m,a^{m}\eta) &= (e,\eta)_{a^{-m}}
\\ &= |a^m| - (a^{-m},\eta)_{e}
\\ &\ge |a^m| - (a^{-m},\gamma)_{e} - 2c_2
\ge |a^m| - c_1 - 2\delta - 2c_2
\ge c_1+3\delta + r,
\end{align*}
if we choose \(m\) large enough depending only on \(c_1,c_2\) and \(\delta\).

We conclude the proof using the hyperbolicity property since
\[
(a,a^m\eta) \ge \min\lbrace (a,a^m), (a^m,a^m\eta) \rbrace -\delta \ge c_1 + 2\delta + r,
\]
so \(a^m\eta \in V(a)\) as required.
\end{proof}

We define
\[
T = \lbrace \gamma: \gamma \in V(a)\text{ and }\gamma^{-1} \in V(a^{-1})\rbrace.
\]
With \(m\) given by the previous lemma, for each \(\gamma \in \Gamma\) let
\[
L(\gamma) =\begin{cases}a^{m+1} &\text{if } \gamma \notin V(a^{-1}),\\ a^{m+1}b &\text{otherwise,}\end{cases}\quand
R(\gamma) =\begin{cases}a &\text{if } (L(\gamma)\gamma)^{-1} \notin V(a),\\ ba &\text{otherwise.}\end{cases}
\]

\begin{lemma}
One has \(L(\gamma)\gamma R(\gamma) \in T\) for all \(\gamma \in \Gamma\).
\end{lemma}
\begin{proof}
Since \(\dist(L(\gamma)\gamma, L(\gamma)\gamma R(\gamma)) \le 2c_2\) this follows immediately from Lem\-ma~\ref{mlemma}.
\end{proof}

We now let \(S\) be the semigroup generated by \(T\).

\begin{lemma}
The semigroup \(S\) generates \(\Gamma\) as a group.
\end{lemma}
\begin{proof}
From the definition \(a \in S\) and, using Lemma \ref{borrowedlemma}, \(aba \in S\). Therefore \(L(\gamma),R(\gamma)\) belong to the group generated by \(S\) for all \(\gamma\).
The equation
\[
\gamma = L(\gamma)^{-1}\left(L(\gamma)\gamma R(\gamma) \right)R(\gamma)^{-1}
\]
now implies that \(S\) generates \(\Gamma\) as a group, as required.
\end{proof}

As a first step towards showing that \(S\) is quasi-geodesic we prove the following.

\begin{lemma}\label{boundeddistanelemma}
For any \(n\) and \(\gamma_1,\ldots, \gamma_n \in T\), letting \(\alpha\) be a geodesic joining \(e\) and \(\gamma_1\cdots \gamma_n\), one has\vspace*{-3pt}
\[
|\gamma_1\cdots\gamma_n| \ge |\gamma_1\cdots \gamma_k| + |\gamma_{k+1}\cdots \gamma_n| - r/8,
\]
and \(\dist(\gamma_1\cdots \gamma_k,\alpha) \le r/8\) for all \(k = 0,\ldots, n\).
\end{lemma}
\begin{proof}
Setting \(x_k = \gamma_1\cdots \gamma_k\) for \(k = 0,\ldots, n\), we observe that\vspace*{-3pt}
\[
(x_{k-1},x_{k+1})_{x_k} = (\gamma_{k}^{-1},\gamma_{k+1}) \le c_1 + 2\delta,
\]
and \(\dist(x_k,x_{k+1}) = |\gamma_{k+1}| \ge r\) by Lemma \ref{borrowedlemma}.
Hence the sequence is a \((c_1+\delta,r)\)-chain as in \cite{gouezel}.
We observe that\vspace*{-3pt}
\[
r > 2(c_1 + 2\delta) +4\delta + 1,
\]
and therefore by \cite[Lem.\,3.8]{gouezel} we have
\[
(x_0,x_n)_{x_k} \le c_1 + 4\delta,
\]
for all \(k\).
This implies that\vspace*{-3pt}
\[
\dist(x_k,\alpha) \le c_1 + 5\delta < r/8.
\]
Also since \((x_0,x_n)_{x_k} = (x_k^{-1},x_k^{-1}x_n)_e\) we get by definition
\[
\frac{|\gamma_1\cdots \gamma_k| + |\gamma_{k+1}\cdots \gamma_n| - |\gamma_1\cdots \gamma_n|}{2} \le c_1 + 4\delta \le r/16,
\]
which yields the required lower bound for \(|\gamma_1\cdots \gamma_n|\).
\end{proof}

We now conclude the proof of Proposition \ref{semigrouplemma}.

\begin{lemma}
The semigroup \(S\) is quasi-geodesic.
\end{lemma}
\begin{proof}
Fix \(n\) and \(\gamma_1,\ldots, \gamma_n \in T\), and let \(\alpha\) be a geodesic with \(\alpha_0 = e\) and \(\alpha_{N_n} = \gamma_1\cdots \gamma_n\) with \(N_n = |\gamma_1\cdots \gamma_n|\).
Since the word distance is integer valued, setting \(C = \lfloor r/8 \rfloor\) we have by Lemma \ref{boundeddistanelemma} that there exist \(N_0 = 0, N_1, \ldots,N_n = |\gamma_1\cdots \gamma_n|\) such that\vspace*{-3pt}
\[
\dist(\gamma_1\cdots \gamma_k, \alpha_{N_k}) \le C \le r/8,
\]
for all \(k\).

We also obtain from Lemma \ref{boundeddistanelemma} applied to \(\gamma_1,\ldots, \gamma_{k+1}\) and Lemma \ref{borrowedlemma} that
\begin{align*}
N_{k+1} - N_k &\ge |\gamma_1\cdots \gamma_{k+1}|-r/8 - |\gamma_1\cdots \gamma_k| -r/8
\\ &\ge |\gamma_{k+1}| - 3r/8 \ge 5r/8 \ge 5C,
\end{align*}
and in particular \(N_{k} \le N_{k+1}\) for \(k = 0,\ldots, n-1\).

We now construct a path \(\beta\) which we claim to be quasi-geodesic (with constants independent of \(n\) and \(\gamma_1,\ldots,\gamma_n\)) by concatenating segments of \(\alpha\) with a paths to and from each \(\gamma_1\cdots \gamma_k\) (which will have length at most \(2C\) as seen above).
To be precise, we choose \(\beta:[0,N_n + 2(n-1)C] \cap \Z \to \Gamma\) such that
\begin{enumerate}
\item For each \(k = 0,\ldots,n-1\) one has \(\beta_{m+N_k+2kC} = \alpha_{m+N_k}\) for all \(m \in [0,N_{k+1}-\nobreak N_k] \cap \Z\).
\item \(\beta_{N_k + 2(k-1)C + C} = \gamma_1\cdots \gamma_k\) for \(k = 1,\ldots, n-1\).
\item \(\dist(\beta_m, \beta_{m+1})\le 1\) for all \(m\).
\end{enumerate}
In view of the third property above we have
\(\dist(\beta_l,\beta_m) \le |l-m|\).
For the lower bound, given \(0 \le l < m \le N_n + 2(n-1)C\) we write
\[
\dist(\beta_l,\beta_m) = \sum_{k = 0}^{m-l}\dist(\beta_{l+k},\beta_{l+k+1}) \ge \sum_{k = 0}^{m-l}f(k),
\]
where we set \(f(k) = 1\) if \(k \in [N_i+2iC,N_{i+1}+2iC]\) for some \(i\) and \(f(k) = -1\) otherwise.
Since the sequence \(f(0),\ldots,f(l-m)\) consists of subsequences of runs of at most \(2C\) consecutive times the value \(-1\), with runs of at least \(5C\) times the value~\(1\) in between, we~obtain
\[
\dist(\beta_l,\beta_m) \ge 3(|m-l| - 4C)/7 - 4C,
\]
which establishes that \(S\) is \((7/3, 6C)\)-quasi-geodesic.
\end{proof}

